\begin{figure}[htbp]
	\centering
	\begin{tikzpicture}[font=\small]
		% ---------------------------------------------------------------
		% (a) trave inclinata dell'angolo alpha
		% ---------------------------------------------------------------
		\begin{scope}
			\def\a{30}
			% riferimento globale
			\draw[trvasse] (0,0) -- (0.9,0) node[right] {$x^{\mathrm{g}}$};
			\draw[trvasse] (0,0) -- (0,0.9) node[above] {$y^{\mathrm{g}}$};
			\draw[line width=.5pt] (0,0) circle (0.09);
			\fill (0,0) circle (0.9pt);
			\node[below left=-1pt] at (0,0) {$z^{\mathrm{g}}$};
			% trave
			\coordinate (A) at (1.5,0.35);
			\coordinate (B) at ($(A)+(\a:3.6)$);
			\draw[trvtrave] (A) -- (B);
			\draw[trvfibra] ($(A)+(\a-90:0.12)$) -- ($(B)+(\a-90:0.12)$);
			\pic at (A) {trvpunto};
			\node[above left=1pt] at (A) {$A$};
			\node[above left=1pt] at (B) {$B$};
			% riferimento locale (in B)
			\draw[trvasse] (B) -- ++(\a:0.85) node[above] {$z$};
			\draw[trvasse] (B) -- ++(\a-90:0.85) node[right] {$y$};
			\draw[line width=.5pt, fill=white] (B) circle (0.09);
			\fill (B) circle (0.9pt);
			\node[right=3pt] at (B) {$x$};
			% angolo alpha
			\draw[trvcostruzione] (A) -- ++(1.6,0);
			\draw[-{Stealth[length=1.3mm,width=1.0mm]},line width=.5pt]
				($(A)+(1.1,0)$) arc[start angle=0, end angle=\a, radius=1.1];
			\node at ($(A)+(\a/2:1.35)$) {$\alpha$};
			\node at (2.7,-1.2) {(a)};
		\end{scope}
		% ---------------------------------------------------------------
		% (b) casi alpha = 0, pi/2, pi, 3pi/2
		% ---------------------------------------------------------------
		\begin{scope}[xshift=6.6cm, yshift=0.2cm]
			\foreach \a/\mx/\my in {0/0.6/1.55, 180/0.6/0.0, 90/3.0/1.25, 270/3.0/-0.15}{
				\pgfmathtruncatemacro{\anz}{mod(\a+180,360)}
				\pgfmathtruncatemacro{\any}{mod(\a+90,360)}
				\coordinate (P) at ($(\mx,\my)+(\a+180:0.55)$);
				\coordinate (Q) at ($(\mx,\my)+(\a:0.55)$);
				\draw[trvtrave] (P) -- (Q);
				\draw[trvfibra] ($(P)+(\a-90:0.11)$) -- ($(Q)+(\a-90:0.11)$);
				\draw[trvasse] (Q) -- ++(\a:0.45) node[font=\footnotesize, anchor=\anz] {$z$};
				\draw[trvasse] (Q) -- ++(\a-90:0.45) node[font=\footnotesize, anchor=\any] {$y$};
			}
			\node[trvnota] at (0.45,2.0) {$\alpha = 0$};
			\node[trvnota] at (0.75,-0.45) {$\alpha = \pi$};
			\node[trvnota, left] at (2.85,1.05) {$\alpha = \tfrac{\pi}{2}$};
			\node[trvnota, right] at (3.15,0.05) {$\alpha = \tfrac{3\pi}{2}$};
			\node at (1.8,-1.4) {(b)};
		\end{scope}
	\end{tikzpicture}
	\caption{Riferimento locale della trave.
		(a)~L'asse $z$ è diretto da $A$ verso $B$ e forma
		l'angolo $\alpha$ con $x^{\mathrm{g}}$; l'asse $y$
		punta verso la fibra inferiore, indicata dal
		tratteggio; l'asse $x$, uscente dal foglio,
		coincide con $z^{\mathrm{g}}$.
		(b)~Riferimento locale per $\alpha = 0$, $\pi/2$,
		$\pi$, $3\pi/2$.}
	\label{fig:riferimento_intrinseco}
\end{figure}
